\clearpage
\VerdierPhysicalPageAnchor{0184}%
\begingroup\small\setlength{\parskip}{0pt}\setlength{\emergencystretch}{1.5em}
\setlength{\abovedisplayskip}{4pt plus 2pt minus 2pt}%
\setlength{\belowdisplayskip}{4pt plus 2pt minus 2pt}%

\VerdierSemanticAnchor{ch3d-1-3-4-c}
\noindent\textit{c) $D(\mathcal A)$ の任意の完全三角は、
$S=0\longrightarrow X\xrightarrow{u}Y\xrightarrow{v}Z\longrightarrow0$
を $\mathcal A$ の完全列として、
$(X,Y,Z,[u],[v],\delta(S))$ の形の完全三角と同型である。}

\VerdierSemanticAnchor{ch3d-1-3-4-proof}
a) と b) の証明は読者に委ねる。c) を証明しよう。（第 II 章，2.2.6）により、
$D(\mathcal A)$ の任意の完全三角は、$K(\mathcal A)$ のある完全三角の
$Q\colon K(\mathcal A)\longrightarrow D(\mathcal A)$ による像と同型である。
（第 I 章，3.3.1）および公理 TRII により、$K(\mathcal A)$ の任意の完全三角は
次の形の三角と同型である：

\VerdierSemanticAnchor{ch3d-1-3-4-cone-triangle}
\[
Y\xrightarrow{\dot{p}(u)}c(\dot u)
\xrightarrow{\dot{q}(u)}X[1]
\xrightarrow{-\dot u[1]}Y[1].
\]
したがって、$D(\mathcal A)$ の任意の完全三角は次の形の三角と同型である：

\VerdierSemanticAnchor{ch3d-1-3-4-derived-triangle}
\[
(Y,c(\dot u),X[1],[p(u)],[q(u)],-[u][1]).
\]
ここで次に注意する：

\VerdierSemanticAnchor{ch3d-1-3-4-cone-exact}
\[
S=0\longrightarrow Y\xrightarrow{p(u)}c(u)
\xrightarrow{q(u)}X[1]\longrightarrow0
\]
は $\operatorname{comp}(\mathcal A)$ の完全列である。したがって
（第 II 章，1.2.4）により、三角
\[
(Y,c(\dot u),X[1],[p(u)],[q(u)],-[u][1])
\]
から次の三角への同型が存在する：
\[
(Y,c(\dot u),X[1],[p(u)],[q(u)],\delta(S)).
\]

\VerdierSemanticAnchor{ch3d-1-3-5}
\medskip\noindent\textbf{命題 1.3.5.}\quad
\textit{a) 擬同型を $C$ の同型に送る任意の関手
$F\colon\operatorname{comp}(\mathcal A)\longrightarrow C$ は、関手
$\operatorname{comp}(\mathcal A)\longrightarrow D(\mathcal A)$ を介して
一意に分解する。}

\noindent\textit{b) $\operatorname{comp}(\mathcal A)$ からアーベル圏 $B$ への
半完全関手であって、擬同型を同型に送るものは、標準 $\delta$-関手
$\operatorname{comp}(\mathcal A)\longrightarrow D(\mathcal A)$ と
$D(\mathcal A)$ から $B$ へのコホモロジー的関手を通じて一意に分解する。}

\VerdierSemanticAnchor{ch3d-1-3-5-proof}
主張 b) は、主張 a) と (1.3.4) による $D(\mathcal A)$ の完全三角の記述から
従う。主張 a) は、$D(\mathcal A)$ の定義（第 II 章，2.2.6）と次の補題から
従う。

\VerdierSemanticAnchor{ch3d-1-3-6}
\medskip\noindent\textbf{補題 1.3.6.}\quad
\textit{$\mathcal A$ を加法圏（必ずしもアーベル圏とは限らない）とし、
$F$ を $\operatorname{comp}(\mathcal A)$ から圏 $C$ への関手とする。
次の条件は同値である：}
\endgroup
